% ======================================================================
% Paper body (v2 revision): the integration-time cost of RFI masking for
% 21 cm intensity-mapping BAO surveys, with CHIME + ATSC DTV as the
% worked example. Numbers regenerate from bao-noise-tolerance:
% scripts/run_forecast.py (figures/tables), scripts/fg_sensitivity.py,
% and are guarded against drift by scripts/check_paper_numbers.py.
% ======================================================================

\section{Introduction}
\label{sec:intro}

Neutral-hydrogen intensity mapping is intended to measure large-scale
structure over wide redshift ranges, including baryon-acoustic-oscillation
(BAO) distances at $z\sim1$--$2.5$ \citep{SeoEisenstein2007,Bull2015}.
CHIME has now detected the cosmological 21\,cm signal in cross-correlation
with galaxy and quasar catalogues \citep{CHIMEDetection2023}, with the
Ly$\alpha$ forest \citep{CHIMELya2024}, and in auto-correlation at
$z\sim1$ \citep{CHIMEAuto2025}. These measurements also make the
treatment of radio-frequency interference (RFI) a practical cosmology
problem rather than only an instrument-characterisation problem.

The 470--608\,MHz part of CHIME's 400--800\,MHz band contains North
American UHF television channels 14--36, each 6\,MHz wide \citep{ATSC53}.
For the 21\,cm line this corresponds to $z=1.33$--$2.02$. Distant
transmitters can appear intermittently even at a protected site,
particularly under enhanced propagation. The usual response is to flag
contaminated time--frequency samples and remove them before mapmaking
\citep[e.g.][]{Offringa2012,CHIMEOverview2022}. The CHIME Overview
forecast treats persistently contaminated frequency ranges as
unavailable, but it does not assign a science cost to the continuum
between a clean channel and a completely dead one
\citep[Fig.~31 and Appendix~A]{CHIMEOverview2022}.

The question that motivated this work is straightforward: can masking in
the DTV band make a BAO measurement impractical even if the detector
performs as intended? A comparison between DTV power and a notional
``BAO power per frame'' does not answer that question. BAO is not a
deterministic waveform that must exceed the interference in an
individual frame. It is a statistical feature of the three-dimensional
power spectrum, recovered by combining many spatial and spectral modes.
The useful comparison is therefore between a masking policy and the
cosmological information that remains after that policy is applied.

This paper isolates the \emph{variance cost of masking}. Starting from a
per-channel masked-time fraction, we (i) map the retained samples to an
effective thermal-noise level, (ii) account for completely removed
bandwidth with a scalar volume-loss approximation, (iii) propagate both
factors through the same Fisher calculation used for the CHIME Overview
forecast, and (iv) invert the result to obtain the observing time
required for a stated BAO target. The two channel-noise reductions used
below are transparent limiting models, not universal bounds on all
mapmaking or power-spectrum estimators. Likewise, the scalar treatment
of excised bandwidth does not include the full spectral window.
Residual RFI that survives the mask is a separate covariance and bias
problem and is not absorbed into an ``effective integration time.''

The software follows the same separation. Simon Foreman's CHIME
adaptation of RadioFisher is retained as a versioned scientific backend
\citep{Foreman2021,Bull2015}; the maintained \textsc{baonoise} package
owns the channel definitions, exact frequency overlaps, masking policy,
Fisher-bank queries, and target calculations. The DTV application is one
worked example of that interface rather than a special case in the
forecasting code. Throughout the paper we emphasize ratios of required
observing time relative to the unmasked calculation. Absolute Fisher
times remain conditional on the idealised foreground, calibration, and
instrument assumptions, whereas the ratios provide a more direct measure
of the additional cost introduced by a specified mask.

\section{Forecast framework}
\label{sec:framework}

We use the Fisher formalism of \citet{Bull2015} in the configuration
adopted for the CHIME collaboration's Overview forecast
\citep[Appendix~A and Table~2]{CHIMEOverview2022}. Foreman's public
\texttt{chime-update} branch implements the as-built four-cylinder
geometry: four $20\,\mathrm{m}\times100\,\mathrm{m}$ cylinders on 22\,m
centres, 256 dual-polarisation feeds per cylinder at 305\,mm pitch, with
78\,m illuminated, and a baseline density $n(u)$ calculated from that
layout \citep{Foreman2021}. The forecast uses a total system temperature
$T_\mathrm{sys}=55$\,K including sky contribution, a survey area of
$31{,}000\,\mathrm{deg}^2$, the 400--800\,MHz band divided into 15
redshift bins, BAO-shift information only, and the idealised foreground
setting $\epsilon_\mathrm{FG}=0$. The fiducial parameters are the Planck
2018 final-release CMB best-fit $\Lambda$CDM values used in that branch
\citep{Planck2018}: $h=0.6732$, $\Omega_m=0.3158$, $\sigma_8=0.812$, and
$\sum m_\nu = 0.06$\,eV.

For a redshift bin centred at $z$, the thermal-noise power entering the
Fisher integrand can be written as
\begin{equation}
\label{eq:noise}
P_N(\mathbf{k};z) =
\frac{T_\mathrm{sys}^2(z)\,S_\mathrm{area}}
     {n_\mathrm{pol}\,t_\mathrm{tot}\,\nu_{21}}\,
\mathcal{I}(\mathbf{k};z).
\end{equation}
Here $\nu_{21}=1420.406$\,MHz, $n_\mathrm{pol}=2$, $t_\mathrm{tot}$ is
the on-sky integration time, and $\mathcal{I}$ contains the baseline
density, beam, and channel-window terms that are independent of
$t_\mathrm{tot}$ \citep[Eq.~A9]{CHIMEOverview2022}. The property used
here is $P_N\propto t_\mathrm{tot}^{-1}$.

Let $F_b^{(0)}(t)$ be the unmasked Fisher matrix for redshift bin $b$ at
integration time $t$. The masking model below returns two scalars for
each bin: a retained-time factor $\bar{w}_b$ and a retained-volume
factor $v_b$. The scenario matrix is then
\begin{equation}
\label{eq:scenario}
F_b(t\,|\,\mathcal{M}) = v_b\,F_b^{(0)}(t\,\bar{w}_b).
\end{equation}
Equation~\eqref{eq:scenario} is exact for the scalar noise-and-volume
model adopted in this paper. It is not an exact representation of
frequency-gap mode coupling or correlated residual contamination. A
bank of $F_b^{(0)}(t)$ matrices therefore evaluates masking scenarios
without rerunning the expensive cosmology calculation, while preserving
the same approximation used by the direct RadioFisher hooks. Direct and
banked evaluations agree to better than 0.2 per cent in $\sigma_A$ for
the tested scenarios.

The BAO observable is the wiggle amplitude $A(z)$, defined through
\begin{equation}
\label{eq:wiggle}
P(k,z) = P_\mathrm{smooth}(k,z)\left[1 + A(z)\,f_\mathrm{bao}(k)\right],
\end{equation}
with fiducial $A_\mathrm{fid}=1$ \citep{SeoEisenstein2007,Bull2015}.
Following the Overview, each redshift bin is analysed independently and
marginalised over $\{A,\alpha_\perp,\alpha_\parallel,b\sigma_8,
f\sigma_8,\sigma_\mathrm{NL}\}$ without external priors. We report the
marginalised per-bin BAO-amplitude uncertainty $\sigma_A(z_b)$; smaller
values correspond to a more precise measurement. When a signal-to-noise
representation is useful, we define
\begin{equation}
\label{eq:snr-bin}
\left(\frac{S}{N}\right)_{\!A,b} =
\frac{A_\mathrm{fid}}{\sigma_A(z_b)} = \frac{1}{\sigma_A(z_b)}.
\end{equation}
For statistically independent redshift bins, the survey-level amplitude
signal-to-noise ratio is
\begin{equation}
\label{eq:snr-survey}
\left(\frac{S}{N}\right)_{\!A,\mathrm{survey}}^{2} =
\sum_b \left(\frac{S}{N}\right)_{\!A,b}^{2} =
\sum_b \sigma_A^{-2}(z_b).
\end{equation}
We use $\sigma_A$ for per-bin plots and $(S/N)_A$ only when information
is combined across bins. We also use the marginalised transverse scale
error $\sigma(D_A)/D_A = \sigma(\alpha_\perp)$ and the volume-distance
combination
\begin{equation}
\label{eq:dv}
\ln D_V = \tfrac{2}{3}\ln\alpha_\perp + \tfrac{1}{3}\ln\alpha_\parallel,
\end{equation}
which is the quantity plotted in Fig.~31 of the Overview. We do not
estimate a 6\,MHz BAO signal-to-noise ratio by inverting that figure.
Instead, the 6\,MHz channel geometry is inserted before the Fisher
matrices are marginalised, so the comparison is made in the forecast
itself.

\section{From a mask table to forecast factors}
\label{sec:mapping}

The model begins with a masked-time fraction $f(\nu)$, so a frequency
slice retains a fraction $w(\nu)=1-f(\nu)$ of the nominal samples. We
assume that flagged samples are removed, that the remaining thermal
noise is diagonal in frequency before the line-of-sight transform, and
that the mask can be summarised by a frequency-dependent duty fraction.
These assumptions exclude residual RFI, time-dependent sky coverage, and
the full covariance generated by an irregular spectral window.

A retained slice has noise variance $\sigma^2(\nu)=\sigma_0^2/w(\nu)$.
RadioFisher assigns one thermal-noise normalisation to each redshift
bin, so this heteroscedastic profile must be reduced to one $\bar{w}_b$.
We use two reductions.

\paragraph{Inverse-variance sample counting.}
If frequency samples are combined with their inverse noise variance,
their information adds in proportion to the number of retained samples.
The effective time factor is
\begin{equation}
\label{eq:wiv}
\bar{w}_{b,\mathrm{IV}} = \langle 1-f \rangle_b, \qquad
t_\mathrm{eff} = t_\mathrm{tot}\,\bar{w}_{b,\mathrm{IV}}.
\end{equation}
This is the fiducial reduction. Under the stated diagonal-noise
assumptions it is the optimal sample-counting result for a mode with
uniform spectral weight.

\paragraph{Unweighted Fourier reduction.}
As a deliberately less adaptive case, suppose the frequency axis is
Fourier transformed without compensating for the channel-dependent
noise. Each radial mode then inherits the arithmetic mean noise power,
giving
\begin{equation}
\label{eq:wf}
\bar{w}_{b,\mathrm{F}} =
\left\langle \frac{1}{1-f} \right\rangle_b^{-1}.
\end{equation}
Jensen's inequality guarantees
$\bar{w}_{b,\mathrm{F}} \le \bar{w}_{b,\mathrm{IV}}$, with equality for
uniform $f$. This statement orders the two reductions constructed here;
it does not prove that every real estimator lies between them.
Calibration errors, correlated noise, nonuniform sky weights, or
additional mode cuts can make a real analysis worse than either scalar
model. For uniform masking, both reductions give $\bar{w}=1-f$, and in
the noise-dominated limit the required integration time increases by
$1/(1-f)$.

\subsection{Exact channel--bin overlap}
\label{sec:overlap}

A 6\,MHz television channel does not generally align with a Fisher
redshift bin. For a bin $z\in[z_b^-,z_b^+]$, the corresponding frequency
interval is
\begin{equation}
\label{eq:bin-freq}
\nu_b^- = \frac{\nu_{21}}{1+z_b^+}, \qquad
\nu_b^+ = \frac{\nu_{21}}{1+z_b^-}.
\end{equation}
ATSC channel $c$ occupies
\begin{equation}
\label{eq:channel-freq}
\nu_c^- = 470 + 6(c-14), \qquad
\nu_c^+ = \nu_c^- + 6 \quad [\mathrm{MHz}].
\end{equation}
Its contribution to bin $b$ is weighted by the exact overlap
\begin{equation}
\label{eq:overlap}
\Delta\nu_{bc} =
\left[\min(\nu_b^+,\nu_c^+) - \max(\nu_b^-,\nu_c^-)\right]_+,
\end{equation}
where $[x]_+=\max(x,0)$. Equations~\eqref{eq:wiv} and \eqref{eq:wf} are
evaluated with these overlap lengths, including the clean part of the
bin. Channels that cross a redshift-bin boundary are split between the
two bins rather than assigned to the nearest centre.

\subsection{Excision as a volume approximation}
\label{sec:excision}

A frequency interval that is removed completely is treated differently
from a retained interval with high noise. If a total bandwidth
$B_{\mathrm{exc},b}$ is excised from a bin of bandwidth $B_b$, we use
\begin{equation}
\label{eq:volume}
v_b = 1 - \frac{B_{\mathrm{exc},b}}{B_b}
\end{equation}
as a first-order loss of physical survey volume. The thermal-noise
factor $\bar{w}_b$ is then computed over the surviving bandwidth only.
This is the same scalar approximation implemented by the RadioFisher
hook; it neglects the detailed $k_\parallel$ window produced by the
gaps.

The excision threshold is not universal. A useful analytic check can be
derived for one contaminated slice occupying a fraction $q$ of a bin.
In the noise-dominated, inverse-variance limit, retaining a slice
masked by fraction $f$ gives
\begin{equation}
\label{eq:excision-ratios}
\begin{aligned}
R_\mathrm{keep} &= \frac{t_\mathrm{keep}}{t_0} = \frac{1}{1-qf},\\[2pt]
R_\mathrm{exc} &= \frac{t_\mathrm{exc}}{t_0} = \frac{1}{\sqrt{1-q}},
\end{aligned}
\end{equation}
where the square root in $R_\mathrm{exc}$ follows because Fisher
information is linear in volume and proportional to $t^2$ in the
noise-dominated limit. The crossover is
\begin{equation}
\label{eq:crossover}
f_\star = \frac{1-\sqrt{1-q}}{q} = \frac{1}{1+\sqrt{1-q}}.
\end{equation}
For small $q$, $f_\star\to0.5$, which motivates the 0.5 default used
for the worked example, but the actual choice should be checked against
the science target and the full Fisher bank.

Channel 30 overlaps 3.84\,MHz of the $z=1.40$--$1.50$ bin, or
$q=0.162$, giving $f_\star=0.522$ in this analytic limit. The adopted
$f=0.97$ is therefore on the excision side of the crossover.
Recovering clean-equivalent depth \emph{within that slice} would
require $1/(1-f)\simeq33$ times as much integration, but this is not a
factor of 33 for the whole bin or survey. In the direct Fisher
calculation, channel-30 excision produces a 0.8 per cent survey-level
penalty and changes the worst-bin 2 per cent $D_A$ target from 0.315 to
0.359 on-sky years. This direct factor, 1.14, exceeds the analytic
$R_\mathrm{exc}=1/\sqrt{1-q}=1.09$ because
Eq.~\eqref{eq:excision-ratios} holds in the noise-dominated limit,
whereas the bin is partially cosmic-variance limited at this target;
volume removed from a partly sample-variance-limited bin cannot be
fully recovered with additional integration time. Retaining the
channel gives 0.373 years under inverse-variance weighting and 1.963
years under the unweighted-Fourier reduction. For this channel and this
target, excision is the lower-cost option under either convention.

\section{Software and reproducibility}
\label{sec:usage}

The software is organised around a narrow boundary between a historical
forecast code and the new masking calculation. The RadioFisher fork
preserves the CHIME implementation and supplies only the Fisher backend
and optional frequency-weight/volume hooks. It is not used as the
public interface for mask construction or scenario analysis, and
unrelated RadioFisher survey paths do not need to be modernised for
this project. The \textsc{baonoise} package contains the channel
definitions, pilot-proxy ingestion, exact overlap calculation,
weighting reductions, excision policy, Fisher-bank interpolation,
marginalisation, and required-time inversion.

The expensive step is the clean Fisher bank. For each redshift bin, the
full matrix is tabulated on a logarithmic integration-time grid.
Building the CHIME bank takes approximately 20 minutes on two cores and
is required only when the instrument, cosmology, redshift bins, or
forecast assumptions change. Masking scenarios then use
Eq.~\eqref{eq:scenario}, reducing a query to interpolation and matrix
marginalisation. A mask table can therefore be evaluated in
milliseconds:

\begin{verbatim}
from baonoise import api, scenarios

fc = api.load()
mask = {17: 0.333,
        30: 0.970,
        31: 0.235}

api.required_time(fc, mask, target=5.0)
api.required_time(fc, mask,
                  target=5.0, zbin=6)
api.tolerance_curve(fc,
                    band=scenarios.DTV_BAND)
\end{verbatim}

The unmasked CHIME calculation is treated as the regression baseline.
The repository writes the per-bin $\sigma(D_V)/D_V$ comparison
alongside every forecast run. The verification suite also compares bank
interpolation with direct Fisher evaluations, checks the expected
noise-dominated and cosmic-variance scaling, and evaluates the same
scenarios through the direct RadioFisher hooks and the banked path.
Unit tests cover channel edges, exact partial overlap, the no-mask
identity, uniform masking, excision versus retained-time accounting,
and agreement between the piecewise overlap calculation and the
frequency-weight hook. Every schema-2 bank records the Bao and
RadioFisher computational-source manifests and digests, cosmological
parameters and physical densities, backend capabilities, redshift edges,
power-spectrum cache identity, experiment settings, and baseline checksum.
The numerical artifact therefore carries the reference bundle needed to
audit the calculation that produced it.

This structure is deliberate. It keeps the original forecast provenance
visible while allowing the maintained package to remain small enough to
test. The same interface can be reused for another experiment after
building a compatible clean Fisher bank. It applies directly to
contaminants that can be represented by a retained-time fraction and/or
excised frequency intervals; it does not imply that every contaminant
has the same residual covariance or bias model.

\section{Worked example: CHIME under DTV masking}
\label{sec:example}

\subsection{Masking scenarios}
\label{sec:scenarios}

We use two scenario families.

\paragraph{Uniform stress tests.}
A common masked fraction $f$ is applied either to the full
470--608\,MHz DTV band or to the complete 400--800\,MHz survey band.
These scenarios isolate the scaling with masked fraction and generate
the tolerance curves (Fig.~\ref{fig:tolerance}).

\paragraph{Pilot-proxy-derived configuration.}
Table~\ref{tab:mask} uses the exposure-weighted per-channel fractions
from the current pilot-proxy survey product \citep{PilotProxy}.
Quarterly fractions are weighted by the corresponding number of valid
frames. Channels 24 and 30 are marked ``refused'' by that calibration
procedure because a clean zero point is not available; the worked
example assigns them $f=0.97$ as an explicit modelling value and
excises them under the default $f_\mathrm{exc}=0.5$ policy. The other
channels retain their measured weighted fractions. This table is the
example used in the numerical results, not a universal model for the
DTV environment at every site or epoch.

The largest effects occur in three redshift regions: channel 30 and the
31--35 cluster near $z\simeq1.4$--$1.5$, channel 24 near
$z\simeq1.65$, and channel 17 near $z\simeq1.9$. We therefore report
both survey-level quantities and targets in the affected bins.

\begin{table}[t]
\centering
\caption{Per-channel masked-time fractions $f_c$ used for the
pilot-proxy-derived worked example. Channel $c$ occupies the 6\,MHz
interval defined by Eq.~\eqref{eq:channel-freq}. Channels marked
``exc.''\ are above the adopted 0.5 threshold and are represented as
volume loss.}
\label{tab:mask}
\small
\begin{tabular}{rl|rl|rl}
\toprule
$c$ & $f_c$ & $c$ & $f_c$ & $c$ & $f_c$ \\
\midrule
14 & 0.012 & 22 & 0.010 & 30 & 0.970 (exc.) \\
15 & 0.010 & 23 & 0.017 & 31 & 0.235 \\
16 & 0.014 & 24 & 0.970 (exc.) & 32 & 0.140 \\
17 & 0.333 & 25 & 0.014 & 33 & 0.101 \\
18 & 0.011 & 26 & 0.040 & 34 & 0.023 \\
19 & 0.012 & 27 & 0.020 & 35 & 0.138 \\
20 & 0.011 & 28 & 0.014 & 36 & 0.012 \\
21 & 0.013 & 29 & 0.013 & & \\
\bottomrule
\end{tabular}
\end{table}

\subsection{Unmasked regression against the Overview forecast}
\label{sec:regression}

The masking calculation is applied only after the unmasked forecast is
validated. At the Overview normalisation of 8,760 on-sky hours, the
clean per-bin volume-distance errors run from 0.47 per cent at
$z=0.85$, through 0.61 per cent at $z=1.45$, to 1.03 per cent at
$z=2.43$. They reproduce the percent-level scale and redshift
dependence of the ideal-CHIME error bars in Fig.~31 of
\citet{CHIMEOverview2022} using the same instrument, cosmology,
binning, information flags, and marginalisation. The corresponding
forecast survey-level BAO-amplitude signal-to-noise ratio is 48.84 at
one on-sky year, and the clean survey-level $(S/N)_A=5$ target requires
208.8 on-sky hours. These machine-readable values are used as the
no-mask regression product; masking is not allowed to change bins
outside its frequency support.

For the pilot-proxy-derived configuration, $\sigma(D_V)/D_V$ at
$z=1.45$ increases from 0.61 to 0.733 per cent, and at $z=1.65$ from
0.706 to 0.852 per cent. Bins outside the DTV redshift range are
unchanged.

\subsection{Results}
\label{sec:results}

Table~\ref{tab:survey} reports the survey-level BAO-amplitude
signal-to-noise ratio after two on-sky years and the integration-time
penalty for a survey-level target of $(S/N)_A=5$. At this low target
the calculation is in the noise-dominated regime, so the same ratios
apply closely to other modest survey-level amplitude targets. The
aggregate $(S/N)_A=5$ target is weak for a CHIME-like Fisher forecast
because the uncontaminated parts of the band carry substantial
information. Per-bin distance and amplitude-S/N targets are therefore
more diagnostic of whether the DTV redshifts remain useful
(Table~\ref{tab:bins}, Fig.~\ref{fig:perbin}).

\begin{table}[t]
\centering
\caption{Survey-level results. $(S/N)_A(2\,\mathrm{yr})$ uses 17,520
on-sky hours. The last column is the time required to reach a
survey-level target of $(S/N)_A=5$, relative to the clean case. Values
in parentheses use the unweighted-Fourier reduction where it differs
from the fiducial inverse-variance result.}
\label{tab:survey}
\small
\setlength{\tabcolsep}{4pt}
\begin{tabular}{lcc}
\toprule
Scenario & $(S/N)_A(2\,\mathrm{yr})$ & $t/t_\mathrm{clean}$ \\
\midrule
No masking & 63.57 & 1.000 \\
Pilot-proxy-derived & 61.96 (61.89) & 1.032 (1.034) \\
25\% of DTV band masked & 60.98 & 1.076 \\
50\% of DTV band masked & 57.88 & 1.152 \\
75\% of DTV band masked & 53.89 & 1.219 \\
97\% of DTV band masked & 49.70 & 1.256 \\
50\% of entire band masked & 48.84 & 2.000 \\
ch30 excised & 63.02 & 1.008 \\
ch30 retained & 63.19 (60.26) & 1.013 (1.064) \\
\bottomrule
\end{tabular}
\end{table}

\begin{table}[t]
\centering
\caption{On-sky time to per-bin targets in the two most affected bins.
One on-sky year is defined as 8,760 accumulated hours, matching the
Overview forecast normalisation. No calendar-duty factor is assumed.}
\label{tab:bins}
\small
\setlength{\tabcolsep}{2.6pt}
\begin{tabular}{llcc}
\toprule
$z$ bin & Scenario & $(S/N)_{A,b}=5$ & $\sigma(D_A)\le2\%$ \\
\midrule
1.40--1.50 & clean & 0.175\,yr & 0.315\,yr \\
1.40--1.50 & pilot-proxy & 0.235\,yr & 0.421\,yr \\
1.40--1.50 & pilot-proxy (F) & 0.237\,yr & 0.424\,yr \\
1.40--1.50 & 50\% DTV masked & 0.350\,yr & 0.629\,yr \\
1.40--1.50 & ch30 retained (F) & 1.091\,yr & 1.963\,yr \\
1.40--1.50 & ch30 excised & 0.201\,yr & 0.359\,yr \\
1.60--1.70 & clean & 0.209\,yr & 0.411\,yr \\
1.60--1.70 & pilot-proxy & 0.286\,yr & 0.547\,yr \\
\bottomrule
\end{tabular}
\end{table}

\begin{figure}[t]
\centering
\includegraphics[width=\columnwidth]{fig2_required_time_vs_masking.pdf}
\caption{On-sky observing time required for three BAO targets as a
function of a uniform masked fraction across 470--608\,MHz. The
isolated markers at the left edge show the pilot-proxy-derived
scenario; their horizontal location is only a plotting convention and
is not an effective uniform mask fraction.}
\label{fig:tolerance}
\end{figure}

The pilot-proxy-derived configuration increases the survey-level time
by 3.2 per cent under the fiducial reduction. In the most affected
bin, the 2 per cent $D_A$ target moves from 0.315 to 0.421 on-sky
years, a factor 1.34. The inverse-variance and unweighted-Fourier
results differ by only 0.5 per cent for that target after the two
near-persistent channels are excised. Most of the cost comes from the
excised bandwidth and the channels with $f\gtrsim0.1$, not from the
channels masked at the one-to-few per cent level.

Uniformly masking half of the DTV band gives a 15.2 per cent
survey-level penalty. In the $z=1.40$--$1.50$ bin, which lies wholly
inside the DTV band, the 2 per cent $D_A$ target doubles from 0.315 to
0.629 on-sky years, as expected for uniform 50 per cent masking in the
noise-dominated limit. The separate 50 per cent whole-band test gives
exactly a factor of two at survey level and serves as an internal
scaling check. The 97 per cent DTV stress test still leaves a moderate
survey-level amplitude penalty because the clean frequencies outside
470--608\,MHz dominate that aggregate target; it should not be read as
evidence that the DTV redshift bins remain precise.

The channel-30 comparison shows why a single survey-level number is not
enough. Excision changes the survey-level time by only 0.8 per cent
but costs 16.2 per cent of the $z=1.40$--$1.50$ bin's scalar volume.
Retaining the channel is only slightly worse under inverse-variance
weighting, but it is a factor 6.24 worse than clean for the bin-level
$D_A$ target under the unweighted-Fourier reduction. The practical
conclusion is narrower than ``always excise persistent RFI'': for
channel 30, at the adopted $f=0.97$ and for the targets considered
here, excision is less expensive than retention.

\begin{figure}[t]
\centering
\includegraphics[width=\columnwidth]{fig4_perbin_significance.pdf}
\caption{Marginalised per-redshift-bin BAO-amplitude uncertainty
$\sigma_A$ after one on-sky year (8,760\,hr), shown as a percentage
because $A_\mathrm{fid}=1$. Larger values correspond to weaker
constraints. Curves are shown for the clean, pilot-proxy-derived, and
50 per cent uniform DTV-band scenarios. The shaded region is the
470--608\,MHz DTV band, and the localized increases near
$z\simeq1.45$, 1.65, and 1.9 follow the exact channel--bin overlaps.}
\label{fig:perbin}
\end{figure}

\subsection{Sensitivity to modelling choices}
\label{sec:robust}

The main numerical conclusions are stable under four checks. First,
once channels 24 and 30 are excised, the two noise reductions give
nearly the same answer for the pilot-proxy-derived scenario: the survey
penalties are 1.032 and 1.034, and the worst-bin $D_A$ times are 0.421
and 0.424 on-sky years. The large difference appears only when a
near-dead slice is deliberately kept in an unweighted transform.

Second, the analysis was repeated with the original five-cylinder CHIME
design study in \citet{Bull2015}, which differs in geometry, noise
model, sky area, cosmology, and parameter treatment. The whole-band 50
per cent test remains exactly a factor of two. The pilot-proxy-derived
worst-bin $D_A$ penalty changes from 1.34 to 1.46, while the
survey-level 50 per cent DTV penalty changes from 1.15 to 1.28 because
the DTV band contributes a different fraction of the total
information. These comparisons support using penalty ratios, but they
do not make the absolute times instrument independent.

Third, in the Bull-configuration bank, increasing the
foreground-residual amplitude from $\epsilon_\mathrm{FG}=10^{-6}$ to
$10^{-5}$ changes the clean survey time by 3.3 per cent and
survey-level masking penalties by at most 0.43 per cent. This is a limited
sensitivity test, not a substitute for a realistic foreground and
calibration covariance.

Fourth, the full CHIME calculation was repeated with the P-ACT-LB
parameters from the Atacama Cosmology Telescope DR6 release
\citep{ACTDR6}. The Planck-2018 and P-ACT-LB banks use the same time grid,
instrument, astrophysical profile, physical-density convention,
power-spectrum generator, and RadioFisher 1.0 backend contract. This matched
comparison separates a change of fiducial cosmology from changes in the
implementation. The clean survey-level five-sigma time changes from
208.8 to 186.9 on-sky hours ($-10.5$ per cent), whereas the masking
penalty changes by at most 0.420 per cent across the scenarios and metrics
tabulated here. Thus the fiducial cosmology materially changes the absolute
normalisation but has negligible effect on the relative masking cost.

All times in the tables and figures are accumulated on-sky hours, with
8,760\,hr labelled as one on-sky year to match the Overview forecast.
We do not infer a transferable calendar-duty factor from a selected
CHIME analysis data set. Night-only selection, calibration quality
cuts, missing days, and RFI masking describe different losses and
cannot be collapsed into one survey efficiency without a separate
accounting model. Calendar time can be reported once an
experiment-specific end-to-end efficiency is specified, but applying
such a factor here would risk double-counting the RFI mask whose cost
we are trying to isolate.

\section{Discussion}
\label{sec:discussion}

\paragraph{What the forecast answers.}
The calculation answers how much additional integration is required
after a specified set of samples has been removed, within the CHIME
Fisher model. It can distinguish a small amount of masking spread
across many channels from a near-persistent transmitter concentrated in
one channel, and it reports the result either for the full survey or
for the redshift bins that actually contain the interference. This is
the appropriate level at which to ask whether masking alone makes a DTV
redshift impractical.

\paragraph{Spectral-window limitations.}
The scalar volume factor in Eq.~\eqref{eq:volume} does not preserve the
exact line-of-sight window. Disconnected frequency support couples
$k_\parallel$ modes, can redistribute foreground leakage, and may
remove more information than the bandwidth fraction alone suggests. An
exact treatment would propagate the mask through the mapmaking or
quadratic-estimator covariance, for example in an $m$-mode framework
\citep{Shaw2014,Shaw2015}. The present calculation should therefore be
read as a first-order tolerance forecast. The approximation is easiest
to defend for small excised fractions and should be revisited when
broad or numerous gaps dominate a bin.

\paragraph{Residual contamination.}
Sub-threshold RFI is not equivalent to missing integration time. A
stochastic, decorrelated residual may add covariance; a coherent,
cyclostationary, or spectrally structured residual can also bias the
recovered BAO parameters. For a detector threshold $\tau$, lowering
the threshold generally masks more data while reducing the residual. A
complete threshold optimisation would minimise a quantity such as
\begin{equation}
\label{eq:mse}
\mathrm{MSE}[\widehat{A};\tau] = \sigma_A^2(\tau) + b_A^2(\tau),
\end{equation}
where this paper supplies the masking contribution to the variance term
and a separate residual model must supply any additional covariance and
the bias $b_A$. The pilot-proxy detector operating curve is the
natural input to that companion calculation.

\paragraph{Software scope.}
The maintained result is the small scenario-to-science-cost interface,
not a wholesale rewrite of RadioFisher. Keeping the CHIME forecast
backend close to Foreman's validated branch reduces the chance that an
unrelated software update changes the reference forecast. New
experiments can use the same mask and target machinery after producing
their own clean Fisher bank. New contaminants can use the same
interface when their first-order effect is sample loss or excised
bandwidth; residual-specific physics requires an additional model
rather than another scalar hook.

\paragraph{Applications beyond this example.}
The inputs are deliberately minimal --- a clean Fisher bank for the
instrument and a per-channel masked-time table for the contaminant ---
so the calculation is not specific to CHIME or to television.
Upcoming and proposed 21\,cm surveys occupy similarly contested
spectrum, including CHORD at 300--1500\,MHz \citep{CHORD2019}, HIRAX
at 400--800\,MHz under a different regulatory environment
\citep{HIRAX2022}, and PUMA at 200--1100\,MHz \citep{PUMA2019}; each
is supported by Fisher-level forecasting from which a clean bank can
be built, after which the mask-to-time machinery applies without
modification. The contaminant need not be broadcast television:
mobile-broadband allocations, aeronautical and radar carriers, or
satellite downlinks enter identically whenever their first-order
effect is flagged time or excised bandwidth. The same conversion
speaks to spectrum management. Protection criteria for radio
astronomy are conventionally stated as tolerable percentages of
observing time affected by interference \citep{ITURRA769}; the
present framework turns such a percentage into the quantity an
observatory actually budgets --- the additional observing time
required to reach a stated science target, survey-wide or in the
specific redshift bins at risk. Finally, for detector designers, the
tolerance curve supplies the variance half of the threshold
optimisation in Eq.~\eqref{eq:mse} for any flagging strategy whose
statistics can be summarised per channel.

\section{Conclusions}
\label{sec:conclusions}

The motivating concern was whether DTV masking by itself could make the
corresponding CHIME BAO redshifts unusable. Within the idealised CHIME
Overview Fisher model and the scalar mask treatment used here, the
answer is no for the pilot-proxy-derived configuration: the
survey-level integration time increases by 3.2 per cent, and the most
affected 2 per cent $D_A$ target increases by a factor 1.34. A uniform
50 per cent mask across the full DTV band increases the survey-level
time by a factor 1.15, although the affected bin-level distance targets
double. The distinction matters because clean frequencies can hide
severe local losses in a survey-wide S/N summary.

For a near-persistent channel, the correct comparison is between
retaining a noisy slice and losing its overlapped volume, not between
instantaneous DTV and BAO powers. Channel 30 overlaps 16.2 per cent of
the relevant redshift bin. At the adopted $f=0.97$, excision is less
expensive than retention for the targets considered here, particularly
for an unweighted spectral transform. The 0.5 excision threshold is a
useful default, not a universal constant; Eq.~\eqref{eq:crossover} and
the Fisher bank provide a direct way to check it for another channel or
target.

These conclusions apply to the variance cost of samples that are
actually removed. They do not show that residual DTV contamination is
harmless, and they do not include the complete mode coupling from
spectral gaps. Combining the masking-cost calculation with a
threshold-dependent residual covariance and bias model is the next step
needed to decide how aggressive an RFI detector should be.

\section*{Data and software availability}

The analysis package and manuscript source distributed with this draft
contain the Fisher banks, the pilot-proxy-derived mask table,
machine-readable result tables, and the scripts used to generate the
figures; the maintained repository is
\url{https://github.com/WVURAIL/bao-noise-tolerance}. The maintained
RadioFisher backend is available from WVURAIL \citep{RadioFisherFork} and
retains Foreman's public CHIME configuration \citep{Foreman2021}. Each
distributed Fisher bank records the scientific source digests and checksums
for its baseline-density file, cosmological configuration, power-spectrum
cache, and experiment settings.
